\documentclass[aspectratio=169,11pt]{beamer}
\usetheme{default}
\usecolortheme{default}

\setbeamertemplate{navigation symbols}{}
\setbeamertemplate{footline}[frame number]

\usepackage{amsmath}
\usepackage{booktabs}
\usepackage{graphicx}

\title{Continuous Analog Computation}
\subtitle{From Acoustic Waves to Chemical Reaction Networks}
\author{Srivijayesh Venugopal}
\date{Cranfield IRP 2025--26}

\begin{document}

% ---------------------------------------------------------------------
\begin{frame}
\titlepage
\end{frame}

% ---------------------------------------------------------------------
\begin{frame}{Why the acoustic direction isn't the right fit}

\begin{itemize}
    \item Operating frequency: 2 kHz in air $\rightarrow$ wavelength $\lambda \approx 17$ cm.
    \item Useful diffractive computing needs many wavelengths; device was only $\sim 3\lambda$ long.
    \item Energy estimate: $E_{\text{op}} \approx 0.24$ nJ, far above a digital MAC.
    \item A different model does not fix size or energy.
    \item Prof. Guo explicitly wants reactions, concentrations, and ion-channel-like dynamics.
\end{itemize}

\vspace{0.5cm}
\centering
\textbf{The substrate is wrong.}

\end{frame}

% ---------------------------------------------------------------------
\begin{frame}{Envelope study}

Estimate physical limits \textit{before} detailed design.

\vspace{0.3cm}
Compare substrates on:
\begin{itemize}
    \item energy per operation,
    \item operating delay,
    \item device size,
    \item experimental realizability.
\end{itemize}

\vspace{0.3cm}
\centering
\textbf{Question: even in the best case, is this worth building?}

\end{frame}

% ---------------------------------------------------------------------
\begin{frame}{Candidate 1: acoustic porous waves}

\begin{columns}
\begin{column}{0.55\textwidth}
energy calculation:
\begin{align*}
    \lambda &= \frac{c_0}{f} = \frac{343}{2000} \approx 0.17\ \text{m} \\[0.3em]
    E_{\text{op}} &\approx I A t \approx 0.24\ \text{nJ} \\[0.3em]
    V &\approx 0.5 \times 0.5 \times 0.01 = 2.5\ \text{L}
\end{align*}
\end{column}
\begin{column}{0.42\textwidth}
\textbf{Pros:}
\begin{itemize}
    \item Mature I/O (speakers, microphones)
\end{itemize}

\vspace{0.3cm}
\textbf{Cons:}
\begin{itemize}
    \item Metre-scale device
    \item Poor energy
    \item Weak nonlinearity
\end{itemize}
\end{column}
\end{columns}

\vspace{0.4cm}
\centering
\textbf{Verdict: simulable, but physically unattractive.}

\end{frame}

% ---------------------------------------------------------------------
\begin{frame}{Candidate 2: chemical reaction networks (CRNs)}

\begin{columns}
\begin{column}{0.55\textwidth}
energy calculation:
\begin{align*}
    E_{\text{event}} &\approx kT \approx 4 \times 10^{-21}\ \text{J} \\[0.3em]
    E_{\text{op}} \text{ for } 10^6 \text{ molecules} &\approx 4\ \text{fJ} \\[0.3em]
    V &\sim \mu\text{L}
\end{align*}
Forward model: mass-action ODEs.
\end{column}
\begin{column}{0.42\textwidth}
\textbf{Pros:}
\begin{itemize}
    \item Tiny energy at small molecule counts
    \item Compact reactors
    \item Natural nonlinear kinetics
    \item Trivial solver
    \item Matches supervisor direction
\end{itemize}

\vspace{0.2cm}
\textbf{Cons:}
\begin{itemize}
    \item Slower than electronics
    \item Rate control and noise matter
\end{itemize}
\end{column}
\end{columns}

\vspace{0.4cm}
\centering
\textbf{Verdict: best trade-offs, aligns with supervisor.}

\end{frame}

% ---------------------------------------------------------------------
\begin{frame}{Candidate 3: reaction-diffusion computing}

\begin{columns}
\begin{column}{0.55\textwidth}
energy calculation:
\begin{align*}
    N &\sim 10^{-3}\ \text{mol/L} \times 10^{-9}\ \text{L} \times 6\times10^{23} \\
      &\approx 6\times10^{11}\ \text{molecules} \\[0.3em]
    E_{\text{op}} &\approx 6\times10^{11} \times 4\times10^{-21}\ \text{J} \\
      &\approx 2\ \text{nJ} \\[0.3em]
    \lambda &\approx \sqrt{D\tau} \approx 10\ \mu\text{m}
\end{align*}
\end{column}
\begin{column}{0.42\textwidth}
\textbf{Pros:}
\begin{itemize}
    \item Naturally spatial
    \item Visual, intuitive dynamics
    \item Good supervisor fit
\end{itemize}

\vspace{0.2cm}
\textbf{Cons:}
\begin{itemize}
    \item Slower and hungrier than well-mixed CRNs
    \item PDE optimisation is fragile
\end{itemize}
\end{column}
\end{columns}

\vspace{0.4cm}
\centering
\textbf{Verdict: strong extension, not the primary model.}

\end{frame}

% ---------------------------------------------------------------------
\begin{frame}{Candidate 4: ionic and electronic alternatives}

\textbf{Ion channels}
\begin{align*}
    E_{\text{ion}} &\approx 10kT \approx 4\times10^{-20}\ \text{J} \\
    E_{\text{op}} \text{ for } 10^6 \text{ ions} &\approx 40\ \text{fJ}
\end{align*}
\begin{itemize}
    \item Energetically best, fast gating.
    \item Hard to reconstitute artificial networks experimentally.
\end{itemize}

\vspace{0.4cm}
\textbf{Organic / iontronic devices (OECTs, porous memristors)}
\begin{align*}
    E_{\text{ideal}} &= \tfrac{1}{2}CV^2 \approx 5\ \text{aJ}
\end{align*}
\begin{itemize}
    \item Very novel, electrical readout.
    \item Coupled electrochemistry + hysteresis; too risky for this timeline.
\end{itemize}

\end{frame}

% ---------------------------------------------------------------------
\begin{frame}{The pivot}

\begin{itemize}
    \item Acoustic waves lose on size, energy, and alignment.
    \item CRNs offer the best combination of feasibility, energy scaling, and supervisor fit.
    \item Reaction-diffusion and ion channels are kept as extensions.
    \item Iontronics is future work.
\end{itemize}

\vspace{0.5cm}
\centering
\Large
\textbf{New direction: a trainable CRN-based Fluid Neural Network.}

\end{frame}

% ---------------------------------------------------------------------
\begin{frame}{CRN Fluid NN: core idea}

State = vector of species concentrations $\mathbf{c}(t)$.

\vspace{0.3cm}
Dynamics = mass-action kinetics:
\begin{equation*}
    \frac{d\mathbf{c}}{dt} = \mathbf{S}\,\mathbf{r}(\mathbf{c}; \mathbf{k})
\end{equation*}

\vspace{0.2cm}
where:
\begin{itemize}
    \item $\mathbf{S}$ = stoichiometry matrix,
    \item $\mathbf{r}$ = reaction-rate vector,
    \item $\mathbf{k}$ = trainable rate constants.
\end{itemize}

\vspace{0.3cm}
Extension: add diffusion $\rightarrow$ reaction-diffusion PDE.

\end{frame}

% ---------------------------------------------------------------------
\begin{frame}{First benchmark: trainable XOR}

Reaction network:
\begin{align*}
    A + B &\rightarrow W \quad \text{(mutual annihilation)} \\
    A &\rightarrow Y \quad \text{(A-only output)} \\
    B &\rightarrow Y \quad \text{(B-only output)} \\
    Y &\rightarrow W \quad \text{(output reset)}
\end{align*}

\vspace{0.3cm}
\begin{itemize}
    \item Inputs: initial concentrations $[A, B]$
    \item Output: final concentration $[Y]$
    \item Trainable parameters: rate constants $k_1, k_2, k_3, k_4$
\end{itemize}

\end{frame}

% ---------------------------------------------------------------------
\begin{frame}{Roadmap to submission}

\begin{table}
\centering
\begin{tabular}{cl}
\toprule
\textbf{Week} & \textbf{Milestone} \\
\midrule
1--2  & Well-mixed CRN learns XOR and a simple classifier \\
3     & Gating: one network, two functions via control species \\
4--5  & 1D/2D reaction-diffusion spatial classifier \\
6     & Orthogonal-species parallelism \\
7     & Digital baseline comparison + physical scaling \\
8     & Write-up, figures, conclusions \\
\bottomrule
\end{tabular}
\end{table}

\end{frame}

% ---------------------------------------------------------------------
\begin{frame}{Conclusion}

\begin{itemize}
    \item Acoustic wave computing is abandoned on first-principles grounds.
    \item CRNs provide a fast, valid, and supervisor-aligned path.
    \item First simulations can run in days.
    \item Target: trainable CRN Fluid NN, with gating, RD, and parallelism extensions.
\end{itemize}

\end{frame}

\end{document}
